\section{Introduction}
\label{sec:intro}

Crop losses attributable to plant disease represent a persistent and
costly challenge for agricultural systems around the world.
The Food and Agriculture Organization of the United Nations estimates
that plant diseases destroy up to 40\% of all food crops
annually~\cite{fao2050}, a figure that translates into food insecurity,
economic hardship and avoidable environmental damage from reactive
pesticide use.
For the billions of people who depend on corn, potato, rice and wheat
as dietary staples, robust early detection combined with a reliable
estimate of how much of the harvest will be affected is among the
most consequential problems that modern data-driven methods can help
address.

Traditional disease surveillance has long depended on trained
agronomists walking through fields and inspecting individual leaves.
The approach is knowledgeable but inherently limited in scale: it is
time-consuming, costly and varies across observers. More critically,
visible symptoms often appear only after an infection has already
spread beyond economically treatable levels~\cite{nawaz2023}.
Convolutional Neural Networks (CNNs) have substantially changed what
is achievable in this space, processing thousands of images per hour
and identifying textural and chromatic disease indicators with
consistent accuracy regardless of time of day or observer
fatigue~\cite{das2024}.
Architectures such as Faster RCNN push this further by localising the
exact tissue regions that are affected, information that is essential
for spatially targeted treatment rather than blanket chemical
application.

Knowing the disease identity is only part of what a farmer
needs. The harder downstream question is quantitative: how much of the
harvest will actually be lost? The answer is shaped by climate
variables including temperature, rainfall and humidity interacting with the
biology of the pathogen and the genetic resilience of the variety.
Ensemble machine learning methods and XGBoost~\cite{chen2016xgb} in
particular, have demonstrated strong capability in capturing these
non-linear relationships between climate and yield~\cite{ml_yield_review}.

\noindent\textbf{Research gap.} Work on plant disease images and work on
yield prediction have largely developed separately: image studies stop at
the disease label or the lesion location, and yield models rarely use
image-derived disease information. Two further gaps affect how far
reported results can be trusted. First, we know of no study that
benchmarks all major CNN families on one crop disease dataset under a
single protocol. Second, when augmented copies of one photograph are
stored as separate files, splits based on filenames can place copies on
both sides of a split and inflate test accuracy; this was the case in our
dataset, and we checked it with image hashing.

\noindent\textbf{Problem statement.} How well can pretrained CNNs and a
pseudo-labelled detector recognise and localise crop diseases when
duplicated images are controlled for, and what can a vision-derived
severity score contribute to yield prediction when image and yield data
are not paired?

\noindent\textbf{Objectives.} We pursue four objectives: (O1) benchmark
pretrained CNNs for crop disease classification under a leakage-audited
protocol; (O2) build a severity-aware detector without manual box
labelling; (O3) show how a vision-derived severity score could feed a
yield model and what the available data can support; and (O4) audit the
data for duplicate leakage and report its effect.

\noindent\textbf{Approach and scope.} We build a three-stage
proof-of-concept pipeline: Stage~1 classifies diseases from leaf images,
Stage~2 localises lesions with a severity-aware Faster RCNN, and Stage~3
passes a severity score to a yield model. Stages~1 and~2 use individual
Bangladeshi leaf images, whereas Stage~3 uses FAO country-year yield
records that are not paired with those images by field, season or
location. Stage~3 is therefore a demonstration of data flow, not a
validated image-to-yield system. The same pipeline could be applied to
matched data, with climate, yield and disease images from the same
fields and seasons, and Section~\ref{sec:stage3} describes what that
would require.

\noindent\textbf{What is new.} Each building block used here (CNN
transfer learning, Faster RCNN, XGBoost) is established; novelty is
not claimed from scale. The contribution lies in four specific
elements: (a) a severity-aware Faster RCNN whose learned head predicts
a continuous lesion-coverage score jointly with boxes and classes,
rather than only class and location; (b) a zero-shot Grounding DINO
annotation pipeline that scales a 280-image manual set to 2,256 images
with no further manual labelling, together with a quantitative
characterisation of that pipeline's failure mode; (c) a proof-of-concept
link from a vision-derived severity signal to a yield
model, which we do not validate; and (d) a duplicate audit based on exact-file and perceptual hashing, which shows that filename-based grouping of augmented copies leaves substantial leakage on this dataset (about 30\% of files are exact duplicates of another file), together with a de-duplicated evaluation protocol that does not require retraining.

\noindent\textbf{Primary contributions:}
\begin{itemize}[noitemsep]
  \item A benchmark of 38 CNN architectures across nine model families on the 31,053-image Bangladeshi dataset (15 classes, including Corn Gray Leaf Spot), evaluated on a de-duplicated test subset ($n=1{,}774$) after an audit showed that filename-based grouping had not prevented duplicate leakage. Seven models (EfficientNetB7, ConvNeXtLarge, EfficientNetB4, ConvNeXtXLarge, EfficientNetB5, EfficientNetV2M and EfficientNetV2S; 93.4 to 93.8\%) are statistically indistinguishable, and the compact EfficientNetV2S is a deployment choice within this group.
  \item A scalable \textbf{zero-shot auto-annotation pipeline} using
    Grounding DINO that grows the annotated training set from 280 to
    2,256 images with no additional manual effort, with annotation
    quality quantified against manual annotations (mean IoU = 0.389).
  \item A \textbf{severity-aware Faster RCNN} model, trained and
    evaluated on a filename-grouped, hash-audited detection split
    (6 of 406 validation images duplicated), reaching mAP@0.5 = 49.3\% on the held-out validation set,
    with a learned severity head achieving
    
    MAE = 0.2206 against Grounding DINO coverage targets.
  \item Grad-CAM heatmaps confirming that ConvNeXtLarge (rank 2 of 38 and statistically tied with the best model) focuses on disease-relevant leaf tissue.
\end{itemize}

\noindent\textbf{Proof of concept (Stage~3).} A dynamic severity score
computed from Faster RCNN box areas and CNN softmax confidence is passed to
an XGBoost yield model. Because the yield records are not paired with the
images, this demonstrates data flow only. A country-held-out audit gives a
negative result: severity is a fixed per-crop constant in training, so it
carries no signal beyond crop identity, and with countries held out XGBoost
(R\textsuperscript{2} = 0.64) does not outperform a crop-mean baseline
(0.64); the R\textsuperscript{2} of 0.97 under a random row split is
inflated by country overlap.
